% Main theorem chapter.

\section{Main theorem}\label{sec:main}

The paper's central claim is a quantitative lower bound on the
mean-squared error of any message-passing graph neural network as
an estimator of the average local clustering coefficient on a
specific distribution over featured graphs. The distribution is
constructed in Section~\ref{sec:entanglement}; every subsequent chapter
either builds a piece of the theorem or applies its consequence to
a case.

\subsection{Statement}

Let $n$ be a positive integer divisible by $12$. Two disjoint finite
families of featured graphs on $n$ vertices, $\mathcal{A}_n$ and
$\mathcal{B}_n$, are defined in Section~\ref{sec:entanglement}. In brief:
$\mathcal{A}_n$ consists of graphs that decompose into $n/6$
monochromatic $(+1)$-triangles and $n/6$ monochromatic
$(-1)$-triangles, so $\Pi_{\clos}$ evaluates to $1$ on every element;
$\mathcal{B}_n$ consists of graphs that decompose into monochromatic
$6$-cycles, so $\Pi_{\clos}$ evaluates to $0$ on every element. Let
$\mathcal{D}_n$ denote the uniform mixture on $\mathcal{A}_n \cup
\mathcal{B}_n$ that assigns total mass $\tfrac{1}{2}$ to each side.

\begin{theorem}[Main theorem]\label{thm:main}
Let $n$ be a positive integer divisible by $12$. Then:
\begin{enumerate}
\item[(i)] \emph{Pointwise MPNN indistinguishability.} For every $M
\in \mpnn$ with continuous update and readout, and every $(\gph_A,
x_A) \in \mathcal{A}_n$ and $(\gph_B, x_B) \in \mathcal{B}_n$,
\[
M(\gph_A, x_A) = M(\gph_B, x_B).
\]
\item[(ii)] \emph{Deterministic closure gap.} $\Pi_{\clos}(\gph_A) =
1$ and $\Pi_{\clos}(\gph_B) = 0$ for every $(\gph_A, x_A) \in
\mathcal{A}_n$ and $(\gph_B, x_B) \in \mathcal{B}_n$.
\item[(iii)] \emph{Quantitative MSE lower bound.} For every $M \in
\mpnn$ with continuous update and readout,
\[
\mathbb{E}_{(\gph, x) \sim \mathcal{D}_n}\!\left[\bigl(M(\gph, x) -
\Pi_{\clos}(\gph, x)\bigr)^2\right] \;\ge\; \frac{1}{4}.
\]
\item[(iv)] \emph{Tightness.} The bound is attained: the constant
MPNN $M \equiv \tfrac{1}{2}$ achieves mean-squared error exactly
$\tfrac{1}{4}$.
\end{enumerate}
The non-MPNN estimator $\hat{\Pi}(\gph, x) = \Pi_{\clos}(\gph, x)$
achieves mean-squared error $0$ on $\mathcal{D}_n$, so the MPNN
class is separated from the optimal estimator by
$\tfrac{1}{4}$ in mean-squared error on this distribution.
\end{theorem}

\subsection{Reading the theorem}

Item (i) is the ensemble-level entanglement result: any MPNN with
continuous update and readout returns the same numerical output on
any pair of featured graphs, one from $\mathcal{A}_n$ and one from
$\mathcal{B}_n$. The output does not depend on which side of the
pair the input is drawn from, so an MPNN cannot use this output to
tell the two families apart.

Item (ii) is the closure gap: every graph in $\mathcal{A}_n$ has
maximum local clustering $\Pi_{\clos} = 1$; every graph in
$\mathcal{B}_n$ has minimum local clustering $\Pi_{\clos} = 0$. The
gap is deterministic, not merely in expectation.

Item (iii) combines (i) and (ii) into a quantitative statement.
Because any MPNN's output is a common value $c$ on all of
$\mathcal{A}_n \cup \mathcal{B}_n$, the estimator's mean-squared
error against $\Pi_{\clos}$ under the uniform mixture $\mathcal{D}_n$
is
\[
\tfrac{1}{2}(c - 1)^2 + \tfrac{1}{2}(c - 0)^2 = c^2 - c + \tfrac{1}{2}
= \bigl(c - \tfrac{1}{2}\bigr)^2 + \tfrac{1}{4}.
\]
The right side is at least $\tfrac{1}{4}$, with equality iff $c =
\tfrac{1}{2}$. This gives (iii) and (iv).

Item (iv) says the bound is tight: a specific MPNN (the constant
$M \equiv \tfrac{1}{2}$, obtainable as a message-passing GNN with
zero-weight aggregation and constant readout) meets the bound. So
the lower bound is not just an inequality; it identifies the
best-case behaviour of any MPNN on this distribution.

The final sentence contrasts MPNNs with the non-MPNN
$\hat{\Pi}$ that returns the correct answer by evaluating
$\Pi_{\clos}$ directly. Zero mean-squared error is achievable
outside the MPNN class; $\tfrac{1}{4}$ is the best achievable inside
the class. The gap of $\tfrac{1}{4}$ is a specific quantitative
statement of how far MPNNs fall short.

\subsection{Robustness under distribution contamination}

The MSE lower bound of Theorem~\ref{thm:main}(iii) survives adding an
arbitrary contamination distribution to $\mathcal{D}_n$.

\begin{corollary}[Robustness]\label{cor:robust}
Fix $n$ divisible by $12$ and $\varepsilon \in [0, 1)$. For any
distribution $\nu$ over featured graphs on $n$ vertices with
features in $\{-1, +1\}$, define the mixture
\[
\tilde{\mathcal{D}}_n^\varepsilon = (1 - \varepsilon)\, \mathcal{D}_n
+ \varepsilon\, \nu.
\]
Then for every $M \in \mpnn$ with continuous update and readout,
\[
\mathbb{E}_{(\gph, x) \sim \tilde{\mathcal{D}}_n^\varepsilon}\!\left[\bigl(M(\gph, x) - \Pi_{\clos}(\gph, x)\bigr)^2\right] \;\ge\; \frac{1 - \varepsilon}{4}.
\]
\end{corollary}

\begin{proof}
The mean-squared error under a mixture decomposes as
\[
\mathbb{E}_{\tilde{\mathcal{D}}_n^\varepsilon}[e^2] = (1 - \varepsilon)\,
\mathbb{E}_{\mathcal{D}_n}[e^2] + \varepsilon \, \mathbb{E}_{\nu}[e^2]
\]
where $e = M - \Pi_{\clos}$. By Theorem~\ref{thm:main}(iii), the first
term is at least $(1-\varepsilon)/4$; the second is non-negative.
\end{proof}

The corollary shows the impossibility is not fragile: any
polynomial adversary that adds an $\varepsilon$-fraction of
arbitrary featured graphs to the training or test distribution
leaves the MSE lower bound at $\Omega(1)$ as long as $\varepsilon <
1$.

\subsection{Structure of the proof}

Item (i) uses the dichotomy of Section~\ref{sec:dichotomy}. Every MPNN
output is a $1$-WL invariant functional of the featured graph
(Proposition~\ref{prop:upper}). The construction of $\mathcal{A}_n$ and
$\mathcal{B}_n$ in Section~\ref{sec:entanglement} arranges for every pair
of samples across the two families to have identical $1$-WL
fingerprints; item (i) follows.

Item (ii) is a direct computation. A vertex in a monochromatic
triangle has local clustering coefficient $1$; a vertex in a
monochromatic $6$-cycle has local clustering coefficient $0$.
Averaging gives (ii).

Items (iii) and (iv) follow by the elementary algebraic argument
sketched above. Corollary~\ref{cor:robust} follows by MSE decomposition on
the contamination mixture.

The pieces are developed in Section~\ref{sec:background} through
Section~\ref{sec:characterization}. The empirical instantiation of the
theorem's consequence for a published estimator (EchoGAE) is in
Section~\ref{sec:emp-reddit}.

\subsection{What the theorem does not say}

The theorem is stated for the average local clustering coefficient
on a specific distribution over structured featured graphs. It has
three known limitations.

First, the ensembles $\mathcal{A}_n$ and $\mathcal{B}_n$ are
structurally uniform. Every connected component is a triangle or a
$6$-cycle, and each component is monochromatic. Real interaction
graphs do not decompose this way. Whether the impossibility extends
to random-graph ensembles that resemble real interaction data is a
separate question addressed in Section~\ref{sec:obstacle}. In the dense
graphon limit, message-passing GNNs are universal for continuous
cut-metric functionals \citep{KerivenPeyre2020}, so the
ensemble-level impossibility does not extend to that regime.

Second, the theorem targets the class of message-passing graph
neural networks with continuous update and readout. Architectures
outside this class (subgraph-count architectures, higher-order
Weisfeiler-Leman variants, spectral architectures with explicit
Laplacian pseudoinverses, Bayesian ideal-point models) are not
subject to the impossibility and can compute the target measures
directly.

Third, the theorem is about $\Pi_{\clos}$; parallel statements hold
for the other closure-based and spectral measures listed in
Proposition~\ref{prop:noncomputable-measures}, with numerical constants
depending on the measure. The choice of $\Pi_{\clos}$ in the
theorem statement is because its definition does not require any
auxiliary partition or opinion vector.
